\subsection{原子、分子、核与粒子}

\begin{结论}{强场 Zeeman：Paschen--Back 区}
$\mu_BB$ 远大于精细结构分裂时 $\mathbf L,\mathbf S$ 解耦，
\[
\重点{\Delta E=\mu_BB\,(m_l+2m_s)}\ (+\ \text{小的 }\zeta m_lm_s\text{ 修正}).
\]
弱场用 $g_J m_J\mu_BB$（Landé），两区之间需对角化 $\mathbf L\cdot\mathbf S+$ Zeeman 项。氢 $2p$ 精细分裂 $4.5\times10^{-5}\,$eV，约 $0.8\,$T 即进入中间区；钠 D 线需几十特斯拉。
\end{结论}

\begin{结论}{双原子分子：转动、振动与同位素位移}
转动 $E_J=BJ(J+1)$，吸收线 $J\to J+1$ 在 $2B(J+1)$，等间距 $2B$。振动含非谐项：
\[
G(v)=\omega_e\bigl(v+\tfrac12\bigr)-\omega_ex_e\bigl(v+\tfrac12\bigr)^2,\qquad \text{Morse 势：}\ \重点{D_e=\frac{\omega_e^2}{4\omega_ex_e}}.
\]
同位素：$\omega_e\propto\mu^{-1/2}$，$B\propto\mu^{-1}$（HCl 中 ${}^{35}$Cl/${}^{37}$Cl 谱线成对）。热布居最大的转动能级 $J_{\max}\approx\sqrt{kT/2B}-\frac12$。
\end{结论}

\begin{结论}{Einstein $A$、$B$ 系数}
以单位角频率能量密度 $u(\omega)$ 计：
\[
\重点{\frac{A_{21}}{B_{21}}=\frac{\hbar\omega^3}{\pi^2c^3}},\qquad g_1B_{12}=g_2B_{21}.
\]
$A\propto\omega^3$：高频（X 射线）跃迁自发辐射极快，难以实现粒子数反转。热平衡下受激/自发之比 $=\frac{1}{\ee^{\hbar\omega/kT}-1}$，可见光在室温下 $\hbar\omega/kT\approx80$--$100$，比值约 $\ee^{-80}$——激光必须靠非平衡泵浦。
\end{结论}

\begin{结论}{谱线展宽：自然、Doppler 与光子流量}
自然展宽 Lorentz 型，半高全宽 $\Delta\nu=\frac{1}{2\pi\tau}$；Doppler 展宽 Gauss 型：
\[
\重点{\Delta\nu_D=\nu_0\sqrt{\frac{8kT\ln2}{mc^2}}}.
\]
室温原子 $\Delta\nu_D/\nu_0\sim10^{-6}$，通常远大于自然线宽（$\sim$MHz）；两者卷积为 Voigt 线型。饱和吸收光谱消除 Doppler 展宽。光子流量 $\dot N=\frac{P\lambda}{hc}$：$1\,$mW、$633\,$nm 约 $3.2\times10^{15}$ 个/秒。
\end{结论}

\begin{结论}{结构因子与消光规则}
$F_{hkl}=\sum_jf_j\ee^{2\pi\ii(hx_j+ky_j+lz_j)}$。立方晶胞：
BCC：\重点{h+k+l\ \text{为偶数}}；FCC：$h,k,l$ 全奇或全偶；金刚石：除 FCC 条件外，全偶时还需 $h+k+l\equiv0\pmod4$。
允许的 $h^2+k^2+l^2$ 序列：简单立方 $1,2,3,4,5,6,8$；BCC $2,4,6,8,10$；FCC $3,4,8,11,12$。比较环半径平方比即可判断晶格类型。
\end{结论}

\begin{结论}{半导体载流子与 pn 结}
质量作用定律 \重点{np=n_i^2}，$n_i=\sqrt{N_cN_v}\,\ee^{-E_g/2kT}$；Si（$E_g=1.12\,$eV）室温 $n_i\approx10^{10}\,\mathrm{cm^{-3}}$。
施主掺杂 $N_D\gg n_i$：$n\approx N_D$，$p=n_i^2/N_D$，Fermi 能级比本征高 $kT\ln\frac{N_D}{n_i}$。理想二极管 $I=I_s(\ee^{qV/kT}-1)$，$kT/q=25.9\,$mV（$300\,$K），电流每增 10 倍电压增约 $60\,$mV。
\end{结论}

\begin{结论}{$\beta$ 衰变与 $\alpha$ 衰变的 $Q$ 值（用原子质量）}
\[
\beta^-:\ Q=[M(Z)-M(Z+1)]c^2,\qquad \beta^+:\ \重点{Q=[M(Z)-M(Z-1)-2m_e]c^2},
\]
电子俘获：$Q=[M(Z)-M(Z-1)]c^2-B_e$（$B_e$ 为被俘获电子的结合能）。
$\beta^+$ 多出的 $2m_ec^2=1.022\,$MeV 来自原子电子数的差。$\alpha$ 衰变（二体）：$T_\alpha=\frac{A-4}{A}Q$，其余为子核反冲。
\end{结论}

\begin{结论}{聚变、裂变与 Lawson 判据}
$\mathrm{D+T\to{}^4He+n}+17.6\,$MeV（中子带走 $14.1\,$MeV，按 $1/m$ 分配）；${}^{235}$U 裂变约 $200\,$MeV。
D--T 点火条件（三重积）：
\[
\重点{n\,T\,\tau_E\gtrsim3\times10^{21}\ \mathrm{keV\,s\,m^{-3}}}\quad(T\approx10\text{--}20\,\mathrm{keV}).
\]
每核子放能：聚变 $\sim3.5\,$MeV，裂变 $\sim0.85\,$MeV；化学反应只有 eV 量级（差 $10^6$ 以上）。
\end{结论}

\begin{结论}{Yukawa 力程与 Sargent 规则}
交换质量 $m$ 的粒子：力程 \重点{R\approx\frac{\hbar}{mc}}，$\pi$ 介子约 $1.4\,$fm，$W$ 玻色子约 $2.5\times10^{-3}\,$fm。
弱衰变宽度 $\Gamma\propto G_F^2m^5$：由 $\mu$ 子寿命推 $\tau$ 轻子：$\tau_\tau\approx\tau_\mu\bigl(\frac{m_\mu}{m_\tau}\bigr)^5\times\mathrm{BR}(\tau\to e\nu\bar\nu)\approx2.2\,\mu\mathrm s\times7.4\times10^{-7}\times0.18\approx2.9\times10^{-13}\,$s，与实测一致。
\end{结论}

\begin{结论}{中微子振荡（两味）}
\[
\重点{P(\nu_\alpha\to\nu_\beta)=\sin^22\theta\,\sin^2\Bigl(\frac{1.27\,\Delta m^2[\mathrm{eV^2}]\,L[\mathrm{km}]}{E[\mathrm{GeV}]}\Bigr)}.
\]
第一振荡极大在 $\frac{L}{E}=\frac{\pi}{2\times1.27\Delta m^2}$：大气 $\Delta m^2\approx2.5\times10^{-3}\,\mathrm{eV^2}$，$E=1\,$GeV 时 $L\approx500\,$km。观测到振荡说明中微子有质量。
\end{结论}

\begin{结论}{守恒律速查与夸克模型磁矩}
强、电磁作用守恒：电荷、重子数、轻子数（各代）、奇异数、宇称；弱作用可破坏奇异数与宇称。判断反应：先查电荷、重子数、轻子数，再看是否需要弱作用（寿命 $\gtrsim10^{-13}\,$s 的典型标志）。
夸克模型 $SU(6)$ 波函数给 \重点{\mu_n/\mu_p=-\tfrac23}（实测 $-0.685$）；$\pi^0\to\gamma\gamma$ 寿命 $\sim10^{-17}\,$s（电磁），$\pi^+\to\mu\nu$ 为 $2.6\times10^{-8}\,$s（弱）。
\end{结论}

\begin{结论}{带电粒子在物质中的能损}
Bethe 公式的标度：\重点{-\frac{\dd E}{\dd x}\propto\frac{z^2}{\beta^2}\,\frac{Z}{A}\rho}（另含对数项），在 $\beta\gamma\approx3$ 附近最小（最小电离粒子），约 $1.5$--$2\,\mathrm{MeV\,cm^2/g}$。
低速时 $1/\beta^2$ 使能损在射程末端急剧增大（Bragg 峰，质子放疗原理）；电子主要经轫致辐射损失能量，辐射长度 $X_0$（水约 $36\,$cm，铅约 $0.56\,$cm）。
\end{结论}

\参考{Foot, \emph{Atomic Physics}；Haken--Wolf, \emph{Molecular Physics and Elements of Quantum Chemistry}；Kittel, \emph{Introduction to Solid State Physics}；Krane, \emph{Introductory Nuclear Physics}；Griffiths, \emph{Introduction to Elementary Particles}；PDG 综述（MIT 8.421 / 8.701 / 8.231）。}
